\%============================================================
\chapter{Bai Hoc: Thien Luong --- Tam Long Tot Lanh La Tai San Quy Nhat (Goodness of Heart --- The Most Precious Wealth)}
\begin{center}
  \textit{\color{truyengold}``Khong co hanh dong tot lanh nao, du nho bao nhieu, la bi lang phi.''}\\[4pt]
  \textit{(No act of kindness, no matter how small, is ever wasted.)}\\[4pt]
  \small\color{truyenblue}--- Co Kim Dong Tay ---
\end{center}
\ngancach

\section{Mother Teresa --- Phuc Vu Nguoi Ngheo La Phuc Vu Chua}
\begin{truyen}{Mother Teresa --- Phuc Vu Nguoi Ngheo La Phuc Vu Chua}{Mother Teresa --- An Do TK 20}
\chuhoa{M}{}M other Teresa den Calcutta nam 1948 va chung kien nhung canh doi nguoi chet doi ngoai pho.\\
\textit{(Mother Teresa arrived in Calcutta in 1948 and witnessed people dying of hunger in the streets.)}

Ba bo truong hoc Loreto --- noi ba day trong an toan --- ra ngoai song voi nhung nguoi ngheo kho nhat.\\
\textit{(She left the Loreto school --- where she taught in safety --- to live with the poorest of the poor.)}

Ba thuong noi voi cac co: 'Moi lan ban cham vao mot thu ve thuong ta noi: Day la Chua trong y phuc kho so'.\\
\textit{(She often told her sisters: 'Every time you touch the wounds of the poor, say: This is Christ in a distressing disguise.')}

Nha cua ba o Calcutta --- Nirdhi House --- chep voi nhung nguoi bi tu bo tren pho, cho chet trong pham gia.\\
\textit{(Her home in Calcutta --- Nirmal Hriday --- filled with people abandoned in the streets, dying with dignity.)}

Ba nhan Giai Nobel Hoa Binh 1979 --- va to chuc cua ba tiep tuc lam viec o 133 quoc gia cho den ngay nay.\\
\textit{(She received the Nobel Peace Prize in 1979 --- and her organization continues working in 133 countries to this day.)}

\end{truyen}
\begin{baihoc}
  Thien luong khong can su ton vinh de tiep tuc --- no tu chay trong tam hon nguoi lam viec vi no la dung.\\
  \textit{(Goodness needs no recognition to continue --- it flows naturally in the heart of one who acts because it is right.)}
\end{baihoc}

\section{Albert Schweitzer --- Bac Si Rung Xanh Chau Phi}
\begin{truyen}{Albert Schweitzer --- Bac Si Rung Xanh Chau Phi}{Albert Schweitzer --- Chau Phi 1913}
\chuhoa{A}{}A lbert Schweitzer o tuoi 30 la mot nha than hoc noi tieng, nhac si dat ten, tac gia nhieu sach ban chay.\\
\textit{(Albert Schweitzer at 30 was a famous theologian, an acclaimed organist, author of several bestselling books.)}

Roi ong doc bai bao keu goi bac si cho vung Chau Phi --- va dat xuong cu quyen sach de di hoc y khoa.\\
\textit{(Then he read an article calling for doctors in Africa --- and set down his books to study medicine.)}

O tuoi 38, ong sang Gabon, Chau Phi, xay mot benh xa bang tre nua giua rung xanh.\\
\textit{(At age 38 he went to Gabon, Africa, and built a hospital from bamboo in the middle of the jungle.)}

Ong lam bac si ky niem o day nam 52 nam --- vo ket hon cung o lai phuc vu cung ong.\\
\textit{(He served as a tropical doctor there for 52 years --- his wife stayed and served alongside him.)}

Ong noi: 'Toi khong biet phan thuong cua cuoc song cua ban la gi --- nhung toi biet rang nen ban phuc vu nguoi khac, ban se hanh phuc.'\\
\textit{(He said: 'I don't know what your destiny will be --- but I know that if you serve others, you will be happy.')}

\end{truyen}
\begin{baihoc}
  Su nghiep lon nhat cua mot con nguoi khong phai la nhung gi ho dat duoc cho ban than --- ma la nhung gi ho giup nguoi khac dat duoc.\\
  \textit{(A person's greatest achievement is not what they attain for themselves --- but what they help others attain.)}
\end{baihoc}

\section{Dr. Norman Bethune --- Bac Si Quoc Te O Mat Tran Trung Quoc}
\begin{truyen}{Dr. Norman Bethune --- Bac Si Quoc Te O Mat Tran Trung Quoc}{Norman Bethune --- Canada/Trung Hoa 1938}
\chuhoa{D}{}D r. Norman Bethune la bac si phau thuat Nguyen Nhat Chinh tai Canada --- nhung ong to ra suc manh lon nhat khi roi nuoc de lai sau lung.\\
\textit{(Dr. Norman Bethune was a leading surgeon in Canada --- but he showed his greatest strength by leaving it all behind.)}

Nam 1938, o tuoi 48, ong di Trung Quoc de phuc vu cac chien binh Cong san chong Nhat.\\
\textit{(In 1938, at age 48, he went to China to serve Communist resistance fighters against Japan.)}

Ong to chuc he thong phau thuat luu dong --- mang phong mo truc tiep den mat tran.\\
\textit{(He organized a mobile surgery system --- bringing the operating room directly to the front lines.)}

Trong mot ca phau thuat khan cap, ong bi dut tay va nhiem trung --- ong tu choi nghi ngoi cho den khi chiu khong noi nua.\\
\textit{(During an emergency surgery his hand was cut and infected --- he refused to rest until he physically could not continue.)}

Ong mat nam 1939; Mao Trach Dong viet bai 'Phuc Vu Nhan Dan' de tuong niem ong --- ten ong duoc biet khap Trung Quoc.\\
\textit{(He died in 1939; Mao Zedong wrote 'Serve the People' in his memory --- his name became known throughout China.)}

\end{truyen}
\begin{baihoc}
  Thien luong khong biet bien gioi quoc gia --- nguoi tot lanh that su tim den noi ho nhat can thiet, du khoang cach la bao xa.\\
  \textit{(Goodness knows no national borders --- the truly good find their way to where they are most needed, no matter how far.)}
\end{baihoc}

\section{Miep Gies --- Nguoi Giau An No Co Anne Frank}
\begin{truyen}{Miep Gies --- Nguoi Giau An No Co Anne Frank}{Miep Gies --- Ha Lan 1942}
\chuhoa{M}{}M iep Gies la mot thu ky binh thuong lam viec cho Otto Frank o Amsterdam trong The chien II.\\
\textit{(Miep Gies was an ordinary secretary working for Otto Frank in Amsterdam during World War II.)}

Khi gia dinh Frank --- bao gom Anne Frank --- phai an tron, Miep toi chap nhan mang thuc pham den cho ho moi ngay.\\
\textit{(When the Frank family --- including Anne Frank --- had to go into hiding, Miep quietly agreed to bring them food every day.)}

Day la hanh dong toi hinh --- tat ca bi bat co the bi xu ban; nhung Miep khong ngan ngai.\\
\textit{(This was a capital offense --- anyone caught could be shot; but Miep did not hesitate.)}

Khi ho bi phat hien, Miep thu thap ban thao nhat ky cua Anne va giu lai --- voi hy vong se tra lai co be sau chien tranh.\\
\textit{(When they were discovered, Miep gathered Anne's diary manuscript and kept it --- hoping to return it to the girl after the war.)}

Anne chet trong trai tap trung, nhung nhat ky cua co, duoc Miep giu gin, tro thanh mot trong nhung cuon sach noi tieng nhat the ky XX.\\
\textit{(Anne died in the camps, but her diary, preserved by Miep, became one of the most famous books of the 20th century.)}

\end{truyen}
\begin{baihoc}
  Hanh dong thien luong khong phai luc nao cung cuu song duoc nguoi ban giu --- nhung no co the cho thay the gioi rang nhan vat cua ho xung dang duoc nho den.\\
  \textit{(Acts of goodness don't always save the one you protect --- but they can show the world that person's life was worth remembering.)}
\end{baihoc}

\section{Nicholas Winton --- Nguoi Cu'u 669 Tre Em Khoi Phat Xit}
\begin{truyen}{Nicholas Winton --- Nguoi Cu'u 669 Tre Em Khoi Phat Xit}{Nicholas Winton --- Anh/Czech 1939}
\chuhoa{I}{}I n 1938, khi nguoi co tuoi su ba muoi, Nicholas Winton di nghi He o Prague --- va bat gap nhung gia dinh Do Thai tuyet vong.\\
\textit{(In 1938, when he was just 29, Nicholas Winton went on holiday to Prague --- and encountered desperate Jewish families.)}

Thay vi nghi tiep, ong o lai va to chuc chuyen tau dua tre em Do Thai den Anh truoc khi Duc xam luoc.\\
\textit{(Instead of continuing his holiday, he stayed and organized trains to bring Jewish children to England before the German invasion.)}

Ong khong duoc phep; ong lam van kien gia, thuyet phuc gia dinh Anh cham soc cac chau be, lo tien ve.\\
\textit{(He had no permission; he forged documents, persuaded English families to care for the children, arranged funding.)}

669 tre em duoc cuu; con tau thu 10 --- lon nhat --- bi can lai boi quoc tang Duc xam luoc; nhung tre tren do cu'u khong thoat.\\
\textit{(669 children were saved; the 10th train --- the largest --- was stopped by the German invasion; the children on board did not escape.)}

Ong khong noi gi ve viec lam cua minh suot 50 nam --- cho den khi vo ong tim thay cuon so trong gac mia chua day ten tat ca 669 dua tre.\\
\textit{(He said nothing about what he had done for 50 years --- until his wife found a scrapbook in his attic containing the names of all 669 children.)}

\end{truyen}
\begin{baihoc}
  Nguoi co tam long tot lanh that su khong can ai biet --- va chinh vi vay, khi su that lo lo, no can hon bao gio het.\\
  \textit{(Those with truly good hearts need no one to know --- and precisely because of that, when the truth finally emerges, it moves us all the more.)}
\end{baihoc}

\section{Clara Barton --- Nu Than Ho Nuoc My}
\begin{truyen}{Clara Barton --- Nu Than Ho Nuoc My}{Clara Barton --- Hoa Ky 1861}
\chuhoa{T}{}T rong noi kinh hoang cua Noi chien My, Clara Barton --- mot phu nu khong co dao tao y te chinh quy --- tu minh den chien truong cham soc thuong binh.\\
\textit{(In the horror of the American Civil War, Clara Barton --- a woman with no formal medical training --- independently went to battlefields to care for the wounded.)}

Ba phan phoi thuoc men, bang bong, do an cho thuong binh ca hai phia Bac va Nam.\\
\textit{(She distributed medicine, bandages, and food to wounded soldiers from both the North and South.)}

Sau chien tranh, ba to chuc tim kiem va xac nhan ten cua 22.000 binh si mat tich --- mot viec chua ai lam.\\
\textit{(After the war, she organized searching for and identifying the names of 22,000 missing soldiers --- work no one had done before.)}

Ba sang lap Hoi Ho Thap Do My nam 1881 va phong vai la chu tich dau tien trong 23 nam.\\
\textit{(She founded the American Red Cross in 1881 and served as its first president for 23 years.)}

Ong noi: 'Toi khong biet co gi tot hon trong cuoc song nay ngoai viec nhin thay nguoi khac thoat khoi dau kho.'\\
\textit{(She said: 'I know nothing better in this life than seeing others relieved from suffering.')}

\end{truyen}
\begin{baihoc}
  Lam dieu tot voi nguoi khong quen biet --- va lam dieu do kien tri suot ca cuoc doi --- la cuoc doi dang song nhat.\\
  \textit{(Doing good for strangers --- and doing so persistently for an entire lifetime --- is the most worthwhile life of all.)}
\end{baihoc}

\section{Raoul Wallenberg --- Nha Ngoai Giao Cuu 100.000 Nguoi Do Thai}
\begin{truyen}{Raoul Wallenberg --- Nha Ngoai Giao Cuu 100.000 Nguoi Do Thai}{Raoul Wallenberg --- Hungary 1944}
\chuhoa{R}{}R aoul Wallenberg, mot nha ngoai giao Thuy Dien tre tuoi, duoc cua den Budapest nam 1944 de 'bao ve' nguoi Do Thai.\\
\textit{(Raoul Wallenberg, a young Swedish diplomat, was sent to Budapest in 1944 to 'protect' the Jewish population.)}

Ong khong co quan doi, khong co phap y --- chi co mot quyen 'ho chieu bao ho' cua chinh ong tu soan --- va long qua cam.\\
\textit{(He had no army, no legal authority --- only a 'protective passport' he had invented himself --- and extraordinary courage.)}

Khi tau hoa chua nguoi Do Thai sap roi ga, ong leo len tau, pha cua, tra 'ho chieu' cho cang nhieu nguoi cang tot truoc mat linh SS.\\
\textit{(When a train carrying Jews was about to leave the station, he climbed onto the train, broke through doors, and handed 'passports' to as many as possible in front of SS soldiers.)}

Ong noi voi bon linh SS: 'Nhung nguoi nay duoi su bao ho cua nuoc Thuy Dien --- ban khong duoc dong dau den ho.'\\
\textit{(He told the SS soldiers: 'These people are under Swedish protection --- you may not touch them.')}

Uoc tinh 100.000 nguoi Do Thai duoc cuu --- ong bi Lien Xo bat giu sau Chien tranh va mat tich --- khong ai biet chinh xac ong chet nhu the nao.\\
\textit{(An estimated 100,000 Jews were saved --- he was arrested by the Soviets after the war and disappeared --- no one knows exactly how he died.)}

\end{truyen}
\begin{baihoc}
  Mot nguoi co tam long thien luong va su qua cam co the lam nhung dieu ma ca mot he thong chinh phu khong the lam.\\
  \textit{(One person with a good heart and extraordinary courage can do what an entire government system could not.)}
\end{baihoc}

\section{Nguyen Trai --- Minh Oan Cho Dan Trong Thu Yen Be}
\begin{truyen}{Nguyen Trai --- Minh Oan Cho Dan Trong Thu Yen Be}{Nguyen Trai --- Viet Nam TK 15}
\chuhoa{N}{}N guyen Trai --- nha chinh tri, nha van hoa lon cua Viet Nam the ky XV --- sau khi giup dat nuoc doc lap, khong xin phan thuong gi.\\
\textit{(Nguyen Trai --- the great statesman and cultural figure of 15th century Vietnam --- after helping win independence, asked for no reward.)}

Ong xin nghí ve o an at Bich Dong --- mot ngoi lang yem tinh giua nui rung Ninh Binh.\\
\textit{(He asked to retire to Bich Dong --- a secluded village in the mountains of Ninh Binh.)}

O day ong day hoc cho dan lang, viet thu phap, lam tho --- va viet Quoc Am Thi Tap bang chu Nom.\\
\textit{(There he taught the villagers, practiced calligraphy, wrote poetry --- and composed a poetry collection in Nom script.)}

Ong viet: 'Bui tra ngon ngot dat Cong Thanh --- Truyen doc may tran dua cho ai.'\\
\textit{(He wrote: 'The dust of fame sweetens the earth of my homeland --- pass down these lessons to those who come after.')}

Chinh tam long yeu nuoc va thuong dan cua moi con nguoi binh di --- khong phai vi quyen loi --- la tinh than Nguyen Trai truyen cho hau the.\\
\textit{(It is the love of country and compassion for the common people --- not for personal gain --- that forms the spirit Nguyen Trai passed to posterity.)}

\end{truyen}
\begin{baihoc}
  Thien luong lon nhat khong can phan thuong, khong can tuong niem --- no chi can duoc thuc hanh trong khi ban con dang song.\\
  \textit{(The greatest goodness needs no reward, no memorial --- it only needs to be practiced while you are still alive.)}
\end{baihoc}

\section{Su Tich An Toan Tre Em --- Eglantyne Jebb}
\begin{truyen}{Su Tich An Toan Tre Em --- Eglantyne Jebb}{Eglantyne Jebb --- Anh 1919}
\chuhoa{E}{}E glantyne Jebb la mot phu nu Anh da ra trai truoc toa an vi phat tan to roi co hinh anh tre em Duc dang chet doi sau The chien I.\\
\textit{(Eglantyne Jebb was an English woman who stood trial for distributing leaflets showing starving German children after World War I.)}

Dan Anh luc do van con cam thu Duc --- anh huong tre em Duc dang chet khong khien ai co cam xuc.\\
\textit{(The British public was still bitter against Germany --- images of starving German children moved no one.)}

Quan toa xu phat ba; chinh vi nhan nhuong cua quan toa da sung coi voi ba, da nop tien bao la cho quy cua ba.\\
\textit{(The judge fined her; moved by her composure the judge himself contributed to her fund.)}

Ba sang lap Cuu Te Tre Em --- tien than cua UNICEF ngay nay --- vi tin rang tre em cua moi quoc gia la vo toi.\\
\textit{(She founded Save the Children --- the precursor of today's UNICEF --- believing the children of every nation are innocent.)}

Ba soat thao 'Bua Bao Ve Tre Em' --- co so cua Cong Uoc LHQ ve Quyen Tre Em 60 nam sau do.\\
\textit{(She drafted the 'Declaration of the Rights of the Child' --- the foundation of the UN Convention on the Rights of the Child 60 years later.)}

\end{truyen}
\begin{baihoc}
  Thien luong khong biet thu --- khong phan biet ban hay thu, tre em My hay tre em Duc, nguoi quen hay nguoi la.\\
  \textit{(Goodness knows no boundaries --- it does not distinguish friend or foe, American or German children, the known or the stranger.)}
\end{baihoc}

\section{Phanh Boi Chau --- Long yeu Nuoc La Su Thien Luong Cao Ca}
\begin{truyen}{Phanh Boi Chau --- Long yeu Nuoc La Su Thien Luong Cao Ca}{Phan Boi Chau --- Viet Nam TK 20}
\chuhoa{P}{}P han Boi Chau dat moi su nghiep ca nhan --- hoc vi, danh tieng, cuoc song an lanh --- de di tim duong cuu nuoc.\\
\textit{(Phan Boi Chau sacrificed every personal achievement --- his degree, reputation, comfortable life --- to seek a path to save his country.)}

Ong sang Nhat, sang Trung Quoc, song chet song canh trong hang chuc nam voi hong vong Viet Nam doc lap.\\
\textit{(He went to Japan, to China, lived on the edge of life and death for decades in hope of Vietnamese independence.)}

Bi Phap bat giam va xu khat --- nhieu dong chi cua ong da bo mang --- nhung ong khong he lung lay xuc cam.\\
\textit{(Arrested and sentenced to death by the French --- many of his comrades perished --- but he never wavered in conviction.)}

Ong viet khi trong tu: 'Sinh the noi nay, khong co gi lon hon tinh yeu dan toc --- va khong co gi nho hon ban than minh.'\\
\textit{(He wrote in prison: 'Born in this land, nothing is greater than love of one's people --- and nothing smaller than oneself.')}

Ong duoc tra tu, bi quan che trong Hue --- nhung tinh than dan toc ong nhen len tiep tuc lan ta trong toan quoc.\\
\textit{(He was released and put under house arrest in Hue --- but the national spirit he ignited continued spreading throughout the country.)}

\end{truyen}
\begin{baihoc}
  Long yeu nuoc thien luong --- khong vi quyen loi ca nhan --- la dang cap cao nhat cua thien luong con nguoi co the dat den.\\
  \textit{(Patriotic goodness --- not for personal gain --- is the highest level of human goodness a person can reach.)}
\end{baihoc}
